\section{综合灵敏度与模型评价}

\subsection{灵敏度与稳健性汇总}

问题一在8个搜索与插值情景下，原始时间偏差极差仅为 $3.56\times10^{-9}$ s；四种插值均给出同一报告偏差。问题二在12个配准情景下时间偏差位于50.1694--50.2699 s，过程噪声放缩十倍造成的轨迹RMS变化不超过0.0654 m。问题三改变配准、HAC带宽、区块长度、过程噪声与门控阈值后，偏差“不启用”结论保持不变。问题四三档各500次压力扰动下九项日程均可行。这些结果分别覆盖数值搜索、模型参数、统计推断和任务边界，不能相互替代。

\subsection{交叉验证}

问题一以连续位置目标与700组精确坐标时差交叉验证；问题二以普通最小二乘和Huber稳健目标对照，并用含离群点合成试验核验；问题三以HAC--Wald和移动区块自助法两条不同原理的推断路线复核；问题四以MILP最优状态、严格裕度上界试验及0.001 s连续时间复核共同验证。各问结论均有独立于主算法单一输出的证据。

\subsection{物理目标与优化目标}

问题四的物理主目标是安全约束内尽可能多完成有效任务。第一级MILP直接最大化任务数量并达到九行容量上界；其后只在任务数固定为9时提高最小裕度和总裕度，因此不会以代理目标牺牲物理主目标。最早可行贪心同样完成9项，但最小裕度仅0.1010；正式方案在数量不变时将其提高到0.4241。

\subsection{工程安全裕度}

连续复核得到射击准备窗最小/最大距离为15.653/18.437 m、最大速度1.025 m/s、最大加速度0.811 m/s$^2$；拍照对应为23.257/27.278 m、0.829 m/s和0.855 m/s$^2$。各连续约束中最小百分比安全裕度为拍照距离上界的31.81\%，相邻任务最小间隔0.10 s，高于0.01 s数值间隔。三档各500次扰动均保持可行，但这不等于真实任务成功概率。

\subsection{模型优点}

第一，时间偏差始终和映射公式成对报告，避免方向歧义；固定公共评价网格使不同候选使用同一数据支持。第二，含噪问题直接融合原始异步事件，不把插值点重复当作新信息；Huber配准、门控关闭调参和创新诊断形成可核验闭环。第三，问题三把统计显著性与工程效应分开，并用区块自助法复核，没有把“不拒绝零偏差”夸大为“证明偏差为零”。第四，问题四将完整准备窗纳入可行性定义，舍入后再加密验证，且以模板上界和可行构造共同证明最大任务数为9。

\subsection{模型局限}

问题一假定时钟差固定，无法处理时钟漂移；三次样条在题面无噪声条件下有效，但不适用于含异常点的原始轨迹。问题二、三没有外部真值，所得空间偏差只能解释为两设备相对差，滤波协方差也是模型内不确定性。问题三10 Hz非事件时刻由前一RTS事件状态短时传播得到，最大传播间隔不超过0.208 s，但并非把网格节点纳入统一RTS后的严格桥接分布。问题四的最小裕度和多样性最优性针对压缩候选集；任务互斥、射击目标唯一性又属于显式建模约定。

\subsection{改进方向}

若获得更长记录，可把常数时间差推广为仿射时钟模型，并用分段残差检验漂移。若能取得高精度外部基准，可识别各传感器的绝对偏差并用交叉验证选择过程噪声。问题三可把全部10 Hz网格作为无观测节点并入事件时间轴，再统一执行RTS平滑；偏差自助法也可在每次重抽样中重新估计时间配准，以覆盖联合不确定性。问题四可直接进行连续时间候选细化，并在任务资源、遮挡、转台角速度等工程条件明确后扩充约束。
